\chapter{双原型记忆的类条件结构约束}

\section{模块定位}
全局域对齐能够缩小源域与目标域的整体差异，但并不能保证跨域同类样本彼此靠近，也不能保证不同故障类别之间保持足够间隔。在 TEP 这类 29 类工业时序诊断任务中，若只强调域间混合，模型可能把不同类别的样本压缩到相近区域，从而产生类别混叠甚至负迁移。为此，RCTA 在可靠性门控伪标签的基础上进一步引入双原型记忆，用显式类别中心约束目标域结构。

本章对应第一章提出的挑战三：全局域对齐容易带来类别混叠，难以保证跨域同类样本聚合、异类样本分离。第三章和第四章分别解决“目标域样本如何进入训练”和“哪些目标域样本可信”，本章则进一步解决“可信目标域样本应被组织到怎样的类别结构中”。

\section{源域原型记忆}
为了显式建模类别结构，RCTA 首先为源域维护原型记忆库~\cite{tanwisuth2021prototype}。设学生网络在源域样本上的特征表示为
\begin{equation}
f_i^s=\phi(x_i^s).
\end{equation}
对第 $k$ 类样本，源域小批量原型定义为
\begin{equation}
\tilde{c}_{k}^{s}=
\frac{1}{|\mathcal{B}_{s}^{k}|}
\sum_{i:y_i^s=k} f_i^s,
\end{equation}
并在实现中做单位化处理。跨批次的源域原型记忆采用动量更新：
\begin{equation}
c_{k}^{s}\leftarrow
\mathrm{Norm}\!\left(
\beta c_k^s+(1-\beta)\tilde{c}_k^s
\right),
\end{equation}
其中 $\beta$ 为原型动量系数，本文取 0.9。源域原型的优势在于标签可靠，能够为各故障类别提供稳定锚点；其不足在于只反映源域类别中心，无法完全覆盖目标域工况下的类别形态。

\section{目标域原型记忆}
为弥补源域原型对目标域结构的表达不足，RCTA 同时维护目标域原型记忆。目标域原型只由通过第四章门控集合 $\Omega_t$ 的样本更新，避免大量低可靠伪标签污染类别中心。对第 $k$ 类门控目标样本，其小批量目标原型可写为
\begin{equation}
\tilde{c}_{k}^{t}=
\frac{1}{|\Omega_t^k|}
\sum_{j\in \Omega_t^k} f_j^t,
\end{equation}
其中 $\Omega_t^k=\{j\in\Omega_t|\hat{y}_j^t=k\}$。目标域原型同样采用动量更新：
\begin{equation}
c_{k}^{t}\leftarrow
\mathrm{Norm}\!\left(
\beta c_k^t+(1-\beta)\tilde{c}_k^t
\right).
\end{equation}
这样，目标域原型会随着训练推进逐渐吸收可信目标样本的信息，使类别中心从纯源域锚点过渡为跨域共享锚点。

\section{参考原型构造}
在实际训练中，参考原型并非简单等于源域原型，而是综合当前源域小批量原型、源域记忆原型和目标域记忆原型构成：
\begin{equation}
c_k^{ref}=
\mathrm{Norm}\!\left(
\alpha_1 \tilde{c}_k^s+\alpha_2 c_k^s+\alpha_3 c_k^t
\right),
\end{equation}
其中权重由各类原型是否可用决定。若某一类别在当前批次中没有可靠目标样本，则参考原型主要由源域信息提供；若目标域原型已经稳定，则参考原型将同时吸收目标域类别结构。

这种设计避免了两个极端：只使用源域原型会忽略目标域工况差异，只使用目标域伪标签原型又容易在训练早期受到噪声影响。双原型记忆通过源域稳定性和目标域适应性的组合，为后续类条件约束提供更平衡的几何中心。

\section{原型吸引损失}
在获得参考原型后，RCTA 同时对源域样本和高可靠目标样本施加原型吸引约束：
\begin{equation}
\mathcal{L}_{att}=
\frac{1}{N}\sum_{i}
\left(1-\cos(f_i,c_{y_i}^{ref})\right).
\end{equation}
对于源域样本，标签 $y_i$ 为真实标签；对于目标域样本，标签由门控后保留的伪标签给出。该损失的几何意义是，在归一化特征空间中把同类样本朝各自类别中心拉近。

源域吸引项能够继续稳定已知类别边界，目标域吸引项则把可信目标样本纳入类别结构。二者共同作用，使模型不再只追求“域不可分”，而是进一步追求“跨域同类可聚合”。

\section{原型分离损失}
仅有类内吸引还不够，为避免原型之间过度接近，进一步定义原型分离损失：
\begin{equation}
\mathcal{L}_{sep}=
\frac{1}{|\mathcal{P}|}
\sum_{(k,l)\in\mathcal{P}}
\max\left(0,\cos(c_k^{ref},c_l^{ref})-m\right),
\end{equation}
其中 $\mathcal{P}$ 表示所有有效类别原型对的集合，$m$ 为分离边界。最终的原型约束写为
\begin{equation}
\mathcal{L}_{proto}=
\mathcal{L}_{att}^s+\mathcal{L}_{att}^t+\lambda_{sep}\mathcal{L}_{sep}.
\end{equation}
这样一来，特征空间不仅被要求“源域与目标域更接近”，还被要求“跨域同类更紧、异类更开”，这正是类条件结构学习相对于纯全局对齐的关键补充。

\section{递进验证中的作用}
在第七章的递进实验中，本章模块作为 C 模块叠加到 A+B 之上。对应比较形式为 A+B 与 A+B+C：若 A+B+C 能在目标域准确率、平衡准确率、混淆矩阵主对角线连续性和 t-SNE 类簇分离度等方面取得更好表现，则说明双原型记忆确实缓解了类别混叠问题。该验证顺序与第一章的挑战--创新点对应关系保持一致。

\section{本章小结}
本章将原 RCTA 框架中的双原型记忆与类条件约束拆分为独立章节，分别说明了源域原型、目标域原型、参考原型、原型吸引和原型分离损失的设计。该环节直接回应多类别工业时序迁移中的类别混叠问题，并为完整 RCTA 提供显式类别结构支撑。
